\section{Ejemplo: Problema P0}
El Problema P0 es un problema de juguete que ilustra, con un enunciado sencillo, cómo se espera que se lean los datos de entrada y se escriban las respuestas.

\subsection{Enunciado}
\textbf{Entrada:} La primera línea contiene la cantidad de casos de prueba $T$. Cada una de las siguientes $T$ líneas contiene una lista de $n$ números enteros ($0 < n \leq 1000$), separados por espacio.

\textbf{Salida:} Por cada caso de prueba se espera una línea de la forma \texttt{cp sp}, donde \texttt{cp} es la cantidad de enteros pares leídos y \texttt{sp} es la suma de esos enteros pares.

\subsection{Caso de ejemplo}
\begin{verbatim}
Entrada
3
1 2 3 4 5 6
10 15 20
7 9 11

Salida
3 12
2 30
0 0
\end{verbatim}

\subsection{Ejecución}
Con los casos de prueba anteriores guardados en \texttt{P0.in}:

\begin{verbatim}
java ProblemaP0 < P0.in > P0.out
python ProblemaP0.py < P0.in > P0.out
./ProblemaP0 < P0.in > P0.out
node ProblemaP0.js < P0.in > P0.out
\end{verbatim}
